\section{Chapter IV. Curves}

\ChapterEntry{p-316}{IV, Remark 2.4.1. Absolute Frobenius over the base field}{302}


\Printed{$F\colon X\to X$ is not $k$-linear. On the contrary, we have a commutative diagram}
\Revised{the absolute Frobenius $F\colon X\to X$ need not be $k$-linear. It is $k$-linear when $k=\mathbf F_p$. In general, we have a commutative diagram}

\ChapterEntry{p-317}{IV, Example 2.5.3. Triviality of an \texorpdfstring{\'etale}{etale} covering}{303}


\Printed{if $X$ is isomorphic to a finite disjoint union of copies of $Y$.}
\Revised{if it is isomorphic over $Y$ to a finite disjoint union of identity coverings $Y\to Y$.}

\ChapterEntry{p-317}{IV, Example 2.5.4. Equality in the genus comparison}{303}


\Printed{we are done. By the way, this shows also that equality occurs (for $f$ separable) only if $n=1$, or $g(Y)=1$ and $f$ is unramified.}
\Revised{we are done. By the way, this shows also that equality occurs (for $f$ separable and $g(Y)\geq1$) only if $n=1$, or $g(Y)=1$ and $f$ is unramified.}

\ChapterEntry{p-318}{IV, Example 2.5.5, note. L\"uroth's theorem in dimension two}{304}


\Printed{An analogous result over $k$ algebraically closed is true also in dimension $2$ (V, \Goto{result-V-6.2.1}{6.2.1}).}
\Revised{Over an algebraically closed field $k$, the dimension-$2$ analogue holds for a subfield $L\subseteq k(t,u)$ over which $k(t,u)$ is finite and separable (V, \Goto{result-V-6.2.1}{6.2.1}).}

\ChapterEntry{p-318}{IV, Exercise 2.2(d). The domain of the permutation action}{304}


\Printed{sets of three distinct elements $\beta_1,\beta_2,\beta_3$ of $k$, $\ne0,1$, as follows: reorder the set $0,1,\infty,\beta_1,\beta_2,\beta_3$ according to a given element $\sigma\in\Sigma_6$,}
\Revised{ordered triples $(\beta_1,\beta_2,\beta_3)$ of distinct elements of $k$, $\ne0,1$, as follows: permute the ordered sextuple $(0,1,\infty,\beta_1,\beta_2,\beta_3)$ according to $\sigma\in\Sigma_6$,}

\Printed{triples of distinct elements $\beta_1,\beta_2,\beta_3$ of $k$, $\ne0,1$, modulo}
\Revised{ordered triples of distinct elements $\beta_1,\beta_2,\beta_3$ of $k$, $\ne0,1$, modulo}

\ChapterEntry{p-318}{IV, Exercise 2.3. Excluding plane lines}{304}


\Printed{curve of degree $d$ in $\Proj^2$. For each point $P\in X$, let $T_P(X)$ be the tangent}
\Revised{curve of degree $d$ with $d\geq2$ in $\Proj^2$. For each point $P\in X$, let $T_P(X)$ be the tangent}

\ChapterEntry{p-319}{IV, Exercise 2.3(h). Hyperflexes and the count of bitangents}{305}


\Printed{A plane quartic curve has exactly $28$ bitangents. (This holds even if the curve has a tangent with four-fold contact, in which case the dual curve $X^*$ has a tacnode.)}
\Revised{A plane quartic with $h$ tangents of four-fold contact has $28-h$ bitangents in the sense of \Goto{ex-IV-2.3b}{(b)}, or $28$ counting those tangents as degenerate bitangents. At each four-fold contact line the dual curve has a unibranch cusp of type $(3,4)$, with $\delta=3$.}

\ChapterEntry{p-320}{IV, Exercise 2.7. The objects classified in degree two}{306}


\Printed{finite \'etale morphisms $f\colon X\to Y$ of degree $2$,}
\Revised{finite \'etale covers $f\colon X\to Y$ of degree $2$ up to $Y$-isomorphism, including the split cover $Y\amalg Y\to Y$,}

\Printed{. Thus an}
\Revised{for a connected cover. For the split cover, the cokernel is directly $\mathcal O_Y$. Thus an}

\ChapterEntry{p-321}{IV, Exercise 2.7, note. Cyclic torsors and Picard torsion}{307}


\Printed{\'etale Galois covers of $Y$ with group $\mathbf Z/n\mathbf Z$ are classified by the \'etale cohomology group $H^1_{\mathrm{\acute{e}t}}(Y,\mathbf Z/n\mathbf Z)$, which is equal to the group of $n$-torsion points of $\operatorname{Pic}Y$.}
\Revised{\'etale cohomology group $H^1_{\mathrm{\acute{e}t}}(Y,\mathbf Z/n\mathbf Z)$ classifies $\mathbf Z/n\mathbf Z$-torsors of $Y$, with the action specified and with disconnected sources allowed. A choice of primitive $n$th root of unity in $k$ identifies this group with the $n$-torsion subgroup of $\operatorname{Pic}Y$.}

\ChapterEntry{p-324}{IV, Proposition 3.5. The union of secants}{310}


\Printed{It is a locally closed subset of $\Proj^n$, of dimension $\leq3$, since (at least locally) it is the image of a morphism from $(X\times X-\Delta)\times\Proj^1$ to $\Proj^n$ which sends $\langle P,Q,t\rangle$ to the point $t$ on the secant line through $P$ and $Q$, suitably parametrized.}
\Revised{It is a constructible subset of $\Proj^n$ whose closure has dimension $\leq3$: over $(X\times X)-\Delta$, the points of the secant lines form a $\Proj^1$-bundle with a natural morphism to $\Proj^n$, and the image is $\operatorname{Sec}X$.}

\ChapterEntry{p-325}{IV, Proposition 3.7. Separability and the converse for the nodal projection criterion}{311}


\Printed{Going back to the proof of (II, \Goto{result-II-7.3}{7.3}), condition \Goto{result-IV-3.7-1}{(1)} says that $\varphi$ is one-to-one almost everywhere, hence birational. When $O$ does lie on a secant line, conditions \Goto{result-IV-3.7-2}{(2)}, \Goto{result-IV-3.7-3}{(3)}, \Goto{result-IV-3.7-4}{(4)} tell us that line meets $X$ in exactly two points $P,Q$, it is not tangent to $X$ at either one, and the tangent lines at $P,Q$ are mapped to distinct lines in $\Proj^2$. Hence the image $\varphi(X)$ has a node at that point.}
\Revised{By \Goto{result-IV-3.7-1}{(1)} the projection is generically injective. By \Goto{result-IV-3.7-2}{(2)} its differential is nonzero, so the projection is separable and hence birational. Each one-point fibre has an immersed branch by \Goto{result-IV-3.7-2}{(2)}, so its image point is smooth. At a secant, \Goto{result-IV-3.7-2}{(2)}, \Goto{result-IV-3.7-3}{(3)}, and \Goto{result-IV-3.7-4}{(4)} give exactly two immersed branches with distinct tangent lines, so the image point is a node. Conversely, a birational projection onto a nodal plane curve is its normalization. It is an isomorphism over the smooth points, and over each node it has two immersed branches with distinct tangent lines. There are finitely many nodes, so \Goto{result-IV-3.7-1}{(1)}--\Goto{result-IV-3.7-4}{(4)} follow.}

\ChapterEntry{p-325}{IV, Proposition 3.8(a). Incidence proof for general multisecants}{311}


\Printed{every secant of $X$ is a multisecant}
\Revised{the pairs of distinct points whose secant is a multisecant are dense in $X\times X$}

\Printed{Fix a point $R$ in $X$, and consider the morphism $\psi\colon X-R\to\Proj^2$ induced by projection from $R$. Since every secant is a multisecant, $\psi$ is a many-to-one map. If $\psi$ is inseparable, then for any $P\in X$, the tangent line $L_P$ at $X$ passes through $R$. This gives \Goto{result-IV-3.8b}{(b)} and our conclusion immediately, so we may assume that each such $\psi$ is separable. In that case, let $T$ be a nonsingular point of $\psi(X)$ over which $\psi$ is not ramified. If $P,Q\in\psi^{-1}(T)$, then the tangent lines $L_P,L_Q$ to $X$ are projected into the tangent line $L_T$ to $\psi(X)$ at $T$. So $L_P$ and $L_Q$ are both in the plane spanned by $R$ and $L_T$, hence coplanar.

Thus we have shown that for any $R$, and for almost all $P,Q$ such that $P,Q,R$ are collinear, $L_P$ and $L_Q$ are coplanar. Therefore, there is an open set of $\langle P,Q\rangle$ in $X\times X$ for which $L_P$ and $L_Q$ are coplanar. But the property of $L_P$ and $L_Q$ being coplanar is a closed condition, so we conclude that for all $P,Q\in X$, $L_P$ and $L_Q$ are coplanar. This is \Goto{result-IV-3.8b}{(b)}.}
\Revised{By the density in \Goto{result-IV-3.8a}{(a)}, there is an irreducible component $W$ of the incidence of distinct collinear triples $(P,Q,R)$ dominating one pair of coordinates. Each projection to a pair has finite fibres, since no line is contained in $X$. Thus $\dim W=2$, and $W$ dominates each pair. For a fixed $R$, projection from $R$ extends to a morphism from $X$. If it is inseparable, then every tangent passes through $R$ by closedness, and \Goto{result-IV-3.8b}{(b)} follows. Otherwise the plane image is smooth and the projection unramified outside a finite set of image points. If the common image $T$ of $P,Q\ne R$ is outside that set, then $L_P$ and $L_Q$ are in the plane spanned by $R$ and the tangent at $T$. Thus the triples in $W$ with noncoplanar tangents have finite fibres over $R$, so their locus has dimension at most one. Coplanarity is closed, so it holds throughout $W$, and, by dominance, for every pair $P,Q$. This is \Goto{result-IV-3.8b}{(b)}.}

\ChapterEntry{p-326}{IV, Theorem 3.9. Projection from a point on the curve}{312}


\Printed{Since $A$ lies on every tangent line to $X$, the corresponding morphism from $X$ to $\Proj^2$ is ramified everywhere. So either the image is a point (in which case $X$ is a line), or it is inseparable \Goto{result-IV-2.2}{(2.2)}.}
\Revised{If $A\in X$, the projection is initially defined on $X-\{A\}$ and extends to a morphism from $X$ to $\Proj^2$. Its differential is zero at the generic point, since $A$ lies on every tangent line to $X$. So either the image is a point (in which case $X$ is a line), or the projection is inseparable.}

\ChapterEntry{p-327}{IV, Theorem 3.9. Characteristic in the conic case}{313}


\Printed{For the conic to be strange, we must have $\operatorname{char}k=2$.}
\Revised{Since $y\in K(X)^p$, all the nonzero intersection multiplicities $v_{P_i}(y)$ are divisible by $p$. Their sum is $d=2$, so $p=2$.}

\ChapterEntry{p-328}{IV, Theorem 3.10. The dimension estimate for bad secants}{314}


\Printed{Therefore, by (\Goto{result-IV-3.8}{3.8}), $X$ has a secant which is not a multisecant, and it has a secant without coplanar tangents. Since the same must be true for nearby secants, we see that there is an open subset of $X\times X$ consisting of pairs $\langle P,Q\rangle$ such that the secant line through $P,Q$ is not a multisecant and does not have coplanar tangents. Hence the subset of $X\times X$ consisting of pairs $\langle P,Q\rangle$ where the secant is a multisecant or has coplanar tangents is a proper subset, has dimension $\le1$, and so the union in $\Proj^3$ of the corresponding secant lines has dimension $\le2$.}
\Revised{By the incidence argument in (\Goto{result-IV-3.8}{3.8}), the constructible locus of pairs giving multisecants is not dense in $X\times X$. The coplanar-tangent locus is closed and proper in the open set of distinct pairs. Each locus therefore has dimension at most one. The corresponding secant lines together give a locus of dimension at most two in $\Proj^3$.}

\Printed{For this we consider the morphism $(X\times X-\Delta)\times\Proj^1\to\Proj^3$ (defined at least locally) which sends $\langle P,Q,t\rangle$ to the point $t$ on the secant line through $P$ and $Q$. If the image has dimension $<3$, then we can choose $O$ lying on no secant. If the image has dimension $=3$, then since it is a morphism between two varieties of the same dimension, we can apply (II, Ex.~\Goto{ex-II-3.7}{3.7}), and find there is an open set of points in $\Proj^3$ over which the fibre is finite.}
\Revised{For this we use the morphism to $\Proj^3$ from the secant $\Proj^1$-bundle over $X\times X-\Delta$ of (\Goto{result-IV-3.5}{3.5}). Its source has dimension three. If it is not dominant, an open set of points of $\Proj^3$ lies on no secants. If it is dominant, by (II, Ex.~\Goto{ex-II-3.7}{3.7}) there is a nonempty open set of points over which the fibre is finite.}

\ChapterEntry{p-329}{IV, Exercise 3.6(a)(1). Rational space quartics}{315}


\Printed{$g=0$, in which case $X$ is either the rational normal quartic in $\Proj^4$ (Ex.~\Goto{ex-IV-3.4}{3.4}) or the rational quartic curve in $\Proj^3$ (II, \Goto{result-II-7.8.6}{7.8.6}), or}
\Revised{$g=0$, in which case the linear span of $X$ is either $\Proj^4$, with its rational normal quartic (Ex.~\Goto{ex-IV-3.4}{3.4}), or $\Proj^3$, where $X$ is a smooth rational quartic obtained by an isomorphic projection of the rational normal quartic. Example (II, \Goto{result-II-7.8.6}{7.8.6}) is one such quartic; or}

\ChapterEntry{p-330}{IV, Exercise 3.11(b). Hypotheses in Severi's theorem}{316}


\Printed{surface in $\Proj^5$ for which there is a projection giving a closed immersion}
\Revised{smooth nondegenerate surface in $\Proj^5$ for which there is a projection giving a closed immersion}

\ChapterEntry{p-332}{IV, Lemma 4.2. Fibres of the elliptic involution}{318}


\Printed{Then $\sigma$ interchanges the two points of each fibre of $g$. Hence $\sigma(P)=Q$, and for any $R\in X$, $R+\sigma(R)$ is a fibre of $g$, hence $R+\sigma(R)\in|P+Q|$, i.e., $R+\sigma(R)\sim P+Q$.}
\Revised{Then $\sigma$ exchanges the two points of an unramified fibre and fixes the point of a ramified fibre. In either case the fibre divisor at $g(R)$ is $R+\sigma(R)$. Hence $\sigma(P)=Q$ and $R+\sigma(R)\in|P+Q|$, i.e., $R+\sigma(R)\sim P+Q$.}

\ChapterEntry{p-333}{IV, Proposition 4.6 and Theorem 4.1(b). Rescaling to Legendre form}{319--320}

\PassagePage{319}


\Printed{Now we make a linear change of $x$ to send $a,b$ to $0,1$, so the equation becomes
\[
y^2=x(x-1)(x-\lambda)
\]
as required.}
\Revised{The roots are distinct by nonsingularity. Write $u=(x-a)/(b-a)$ and choose $s\in k^\times$ with $s^2=(b-a)^3$. Setting $v=y/s$ and $\lambda=(c-a)/(b-a)$, we obtain
\[
v^2=u(u-1)(u-\lambda)
\]
as required.}

\PassagePage{320}


\Printed{of variable in $x$, we have $\lambda=\lambda'$.}
\Revised{of variable in $x$ and the corresponding rescaling of $y$ as in (\Goto{result-IV-4.6}{4.6}), we have $\lambda=\lambda'$.}

\ChapterEntry{p-334}{IV, Example 4.6.2 and Corollary 4.7. Cube roots of unity}{320}


\Printed{$\omega^3=1$. Therefore, $j=0$.}
\Revised{$\omega$ is a primitive cube root of unity. Therefore, $j=0$.}

\ChapterEntry{p-335}{IV, Corollary 4.7. Lifts and stabilizer cardinalities}{321}


\Printed{If $\tau=\operatorname{id}$, then either $\sigma=\operatorname{id}$ or $\sigma$ is the automorphism interchanging the sheets of $f$. Thus in any case we have two elements in $G$.

If $\tau\ne\operatorname{id}$, then $\tau$ permutes $\{0,1,\lambda\}$, so $\lambda$ must be equal to one of the other expressions of (\Goto{result-IV-4.5}{4.5}). This can happen only in the following cases:}
\Revised{Conversely, write such an affine symmetry as $\tau(x)=ax+b$. With $F_\lambda(x)=x(x-1)(x-\lambda)$ we have $F_\lambda(ax+b)=a^3F_\lambda(x)$. The two maps $(x,y)\mapsto(ax+b,sy)$, with $s^2=a^3$, are its lifts to $X$ fixing $P_0$. These are all the lifts, since the degree-two extension has only two automorphisms over $\Proj^1$. Thus $|G|=2|\operatorname{Stab}_{\Sigma_3}(\lambda)|$. By (\Goto{result-IV-4.5}{4.5}), the stabilizer is trivial except in the following cases:}

\Printed{$\lambda$ coincides with one other element of its orbit under $\Sigma_3$}
\Revised{the stabilizer has order $2$}

\Printed{and $\operatorname{char}k\ne3$, then}
\Revised{where $\omega$ is a primitive cube root of unity, and $\operatorname{char}k\ne3$, then}

\Printed{$\lambda$ coincides with two other elements of its orbit under $\Sigma_3$}
\Revised{the stabilizer has order $3$}

\Printed{all six elements of the orbit are the same}
\Revised{the stabilizer has order $6$}

\ChapterEntry{p-335}{IV, \S4. Collinearity with multiplicities in the elliptic group law}{321}


\Printed{three points $P,Q,R$ of the image are collinear if and only if $P+Q+R\sim3P_0$.}
\Revised{a line cuts out the divisor $P+Q+R$ if and only if $P+Q+R\sim3P_0$. For three distinct points this is ordinary collinearity. Repeated points are counted with the corresponding intersection multiplicities.}

\ChapterEntry{p-336}{IV, Proposition 4.8. Regularity of the chord construction}{322}


\Printed{Form the equation of the line $L$ joining $P$ and $Q$. This depends on the coordinates of $P$ and $Q$. Now intersect $L$ with $X$. We get a cubic equation in the parameter along $L$, but we already know two of the intersections, so we obtain the coordinates of the third point of intersection $R$ as rational functions in the coordinates of $P$ and $Q$. Since $R=-P-Q$ in the group structure, this shows that the map $(X\times X-\Delta)\to X$ defined by $\langle P,Q\rangle\mapsto-P-Q$ is a morphism.}
\Revised{Locally off the diagonal, choose regular homogeneous representatives $p,q$ of $P,Q$. If $F$ is the cubic equation, then
\[
F(up+vq)=uv(Au+Bv)
\]
with regular coefficients $A,B$. They cannot vanish simultaneously, since the line joining $P$ and $Q$ is not contained in $X$. Thus the residual point $R$, represented by $-Bp+Aq$, varies regularly. Since $R=-P-Q$ in the group structure, this defines a morphism $(X\times X-\Delta)\to X$.}

\ChapterEntry{p-336}{IV, Examples 4.8.1--4.8.3. Torsion points and degrees of multiplication}{322}


\Printed{its kernel is a group isomorphic to $\mathbf Z/n\times\mathbf Z/n$ if $(n,p)=1$, where $p=\operatorname{char}k$, and is isomorphic to $\mathbf Z/p$ or $0$ if $n=p$, depending on the Hasse invariant of $X$.}
\Revised{for $n\ne0$, its kernel on $k$-points is $X[n](k)$, including $P_0$. This group is isomorphic to $\mathbf Z/n\times\mathbf Z/n$ if $\operatorname{char}k=0$ or if $p=\operatorname{char}k>0$ and $(n,p)=1$, and to $\mathbf Z/p$ or $0$ if $n=p>0$, depending on the Hasse invariant of $X$. The differential of $\mu$ at $(P_0,P_0)$ adds the two tangent vectors, since its restriction to each factor is the identity. Hence the differential of $n_X$ at $P_0$ is multiplication by $n$. When this is nonzero, $n_X$ is separable and translations show that it is unramified everywhere. Its degree then equals $|X[n](k)|$.}

\Printed{a point of order $2$ on $X$}
\Revised{in $X[2](k)$}

\Printed{we see that it has degree $4$, and its kernel}
\Revised{it is separable and unramified, so it has degree $4$, and its kernel}

\Printed{is $\mathbf Z/2\times\mathbf Z/2$.}
\Revised{on $k$-points is $\mathbf Z/2\times\mathbf Z/2$.}

\Printed{a point of order $3$ on $X$}
\Revised{in $X[3](k)$}

\Printed{Thus $3_X$ has degree $9$, and its kernel is isomorphic to $\mathbf Z/3\times\mathbf Z/3$.}
\Revised{For $\operatorname{char}k\ne2,3$, it follows that $3_X$ is separable and unramified of degree $9$, with kernel $\mathbf Z/3\times\mathbf Z/3$ on $k$-points.}

\Printed{inflection points of $X$, then the line $PQ$ meets $X$ in a third}
\Revised{distinct inflection points of $X$, then the line $PQ$ meets $X$ in a third}

\Printed{a point of order $3$}
\Revised{in $X[3](k)$}

\ChapterEntry{p-336}{IV, Lemma 4.9. Constant morphisms}{322}


\Printed{If $P+Q=R$ on $X$, then $P+Q\sim R+P_0$ as divisors.}
\Revised{If $f$ is constant, it is the zero homomorphism. Otherwise $f$ is a finite morphism of curves. If $P+Q=R$ on $X$, then $P+Q\sim R+P_0$ as divisors.}

\ChapterEntry{p-337}{IV, \S4. Addition of endomorphisms}{323}


\Printed{$f\times g$$\colon X\to X\times X$ with $\mu$.}
\Revised{$(f,g)$$\colon X\to X\times X$ with $\mu$.}

\ChapterEntry{p-338}{IV, \S4. Pullback of the universal Jacobian family}{324}


\Printed{$f$$\colon X\times T\to X\times J$ induces a homomorphism}
\Revised{$\operatorname{id}_X\times f$$\colon X\times T\to X\times J$ induces a homomorphism}

\ChapterEntry{p-338}{IV, \S4. Unit and inverse identities for a group scheme}{324}


\Printed{A scheme $X$ with a morphism to another scheme $S$ is a \emph{group scheme over $S$} if there is a section $e\colon S\to X$ (the identity) and a morphism $\rho\colon X\to X$ over $S$ (the inverse) and a morphism $\mu\colon X\times X\to X$ over $S$ (the group operation) such that
\begin{enumerate}[label=(\arabic*),leftmargin=2.5em,topsep=3pt,itemsep=1pt]
\item the composition $\mu\circ(\operatorname{id}\times\rho)\colon X\to X$ is equal to the projection $X\to S$ followed by $e$, and
\item the two morphisms $\mu\circ(\mu\times\operatorname{id})$ and $\mu\circ(\operatorname{id}\times\mu)$ from $X\times X\times X\to X$ are the same.
\end{enumerate}}
\Revised{A scheme $X$ with a morphism $\pi\colon X\to S$ is a \emph{group scheme over $S$} if there is a section $e\colon S\to X$ (the identity), a morphism $\rho\colon X\to X$ over $S$ (the inverse), and a morphism $\mu\colon X\times_S X\to X$ over $S$ (the group operation) such that, writing $\varepsilon=e\circ\pi$,
\begin{align*}
\mu\circ(\varepsilon,\operatorname{id}_X)&=\operatorname{id}_X=\mu\circ(\operatorname{id}_X,\varepsilon),\\
\mu\circ(\rho,\operatorname{id}_X)&=\varepsilon=\mu\circ(\operatorname{id}_X,\rho),
\end{align*}
and the two morphisms $\mu\circ(\mu\times_S\operatorname{id}_X)$ and $\mu\circ(\operatorname{id}_X\times_S\mu)$ from $X\times_S X\times_S X$ to $X$ are the same.}

\Printed{the properties (\Goto{cond-IV-group-inverse}{1}), (\Goto{cond-IV-group-assoc}{2}) can be checked on the closed points of $X$. Then (\Goto{cond-IV-group-inverse}{1}) says that $\rho$ gives the inverse of each point, and (\Goto{cond-IV-group-assoc}{2}) says that the group law is associative.}
\Revised{the identity, inverse and associativity laws can be checked on the closed points of the respective products.}

\Printed{properties (\Goto{cond-IV-group-inverse}{1}) and (\Goto{cond-IV-group-assoc}{2})}
\Revised{group identities}

\ChapterEntry{p-339}{IV, Remark 4.10.9. The genus-zero case and translations}{325}


\Printed{If we fix a base point $P_0\in X$, then for any $n\ge1$ there is a}
\Revised{If $g=0$, then $X\cong\Proj^1$ and $\operatorname{Pic}^{\circ}X=0$. By (\Goto{result-IV-4.10.4}{4.10.4}) and (\Goto{result-IV-4.10.7}{4.10.7}), $J$ has just one closed point with zero tangent space, so $J=\operatorname{Spec}k$. Suppose now that $g\ge1$. If we fix a base point $P_0\in X$, then for any $n\ge1$ there is a}

\Printed{Since it is a group scheme, it is a homogeneous space, hence nonsingular everywhere.}
\Revised{Since $J$ is a group scheme, translation by a closed point is an automorphism carrying $0$ to that point. Hence $J$ is regular at every closed point, and therefore everywhere.}

\ChapterEntry{p-340}{IV, Theorem 4.11. The degree-one divisor is a section}{326}


\Printed{which defines an effective Cartier divisor $Z\subseteq X\times T$. By construction, $Z$ intersects each fibre of $p$ in just one point, and in fact one sees easily that the restricted morphism $p\colon Z\to T$ is an isomorphism.}
\Revised{whose restriction to every fibre is a nonzero section of an invertible sheaf of degree one. Its zero scheme $Z$ is a relative effective Cartier divisor, flat over $T$, by the fibrewise criterion, since $X\times T$ is flat over $T$ and the restriction cuts out an effective Cartier divisor on each fibre. The projection $Z\to T$ is proper and quasi-finite, hence finite. The fibres have length one, so $p_*\mathcal O_Z$ is locally free of rank one and its unit generates every residue fibre. Thus $\mathcal O_T\xrightarrow{\sim}p_*\mathcal O_Z$ and $Z\to T$ is an isomorphism.}

\ChapterEntry{p-342}{IV, Theorem 4.13B. Zeros and poles with coincident points}{328}


\Printed{with zeros at $a_1$ and $a_2$, and poles at $b$ and $0$.}
\Revised{with divisor $(a_1)+(a_2)-(b)-(0)$, with coefficients combined at coincident points.}

\ChapterEntry{p-342}{IV, Theorem 4.15B. Domain and fundamental region}{328}


\Printed{complex numbers. Then $J(\tau)=J(\tau')$ if and only if there are integers}
\Revised{nonreal complex numbers. Then $J(\tau)=J(\tau')$ if and only if there are integers}

\Printed{\[
-\tfrac12\le\operatorname{Re}\tau<\tfrac12
\]}
\Revised{\[
\operatorname{Im}\tau>0,\qquad
-\tfrac12\le\operatorname{Re}\tau<\tfrac12
\]}

\ChapterEntry{p-343}{IV, Theorem 4.16 and Corollary 4.17. Torsion and finite multiplication maps}{329}


\Printed{for any $n$, the subgroup of points of order $n$}
\Revised{for any $n\ge1$, the subgroup of points of order dividing $n$}

\Printed{points of order $n$ are represented by $(a/n)+(b/n)\tau$, with $a,b=0,1,\ldots,n-1$.}
\Revised{points of order dividing $n$ are represented by $(a/n)+(b/n)\tau$, with $a,b=0,1,\ldots,n-1$.}

\Printed{morphism multiplication by $n$, $n_X\colon X\to X$, is a finite morphism of degree $n^2$.}
\Revised{morphism multiplication by $n$ for $n\ne0$, $n_X\colon X\to X$, is a finite morphism of degree $n^2$.}

\ChapterEntry{p-343}{IV, Proposition 4.18. Lifting to the universal cover}{329}


\Printed{Hence under the identification of $X$ with $\mathbf C/\Lambda$ it gives a group homomorphism $\bar f$ of $\mathbf C$ to $\mathbf C$, such that $\bar f(\Lambda)\subseteq\Lambda$. On the other hand, since $f$ is a morphism, the induced map $\bar f\colon\mathbf C\to\mathbf C$ is holomorphic. Now expanding $\bar f$ as a power series in a neighborhood of the origin, and expressing the fact that $\bar f(z+w)=\bar f(z)+\bar f(w)$ for any $z,w$ there, we see that $\bar f$ must be just multiplication by some complex number $\alpha$.}
\Revised{Under the identification of $X$ with $\mathbf C/\Lambda$, let $\pi\colon\mathbf C\to\mathbf C/\Lambda$ be the covering projection. The composite $f\circ\pi$ has a unique holomorphic lift $\widetilde f\colon\mathbf C\to\mathbf C$ sending $0$ to $0$. The difference $\widetilde f(z+w)-\widetilde f(z)-\widetilde f(w)$ lies in the discrete lattice $\Lambda$ and is continuous, so it is constant. It is zero at $(0,0)$. Thus $\widetilde f$ is additive and holomorphic, hence multiplication by a complex number $\alpha$. Since $\widetilde f(\Lambda)\subseteq\Lambda$, $\alpha\cdot\Lambda\subseteq\Lambda$.}

\ChapterEntry{p-344}{IV, Theorem 4.19. Trace of complex multiplication}{330}


\Printed{\[
\alpha^2-(a-e)\alpha+(ae-bc)=0,
\]}
\Revised{\[
\alpha^2-(a+e)\alpha+(ae-bc)=0,
\]}

\ChapterEntry{p-345}{IV, Example 4.20.1. Automorphisms fixing the origin}{331}


\Printed{of automorphisms of $X$ has order $4$, so by (\Goto{result-IV-4.7}{4.7}) we must have $j=1728$}
\Revised{$\operatorname{Aut}(X,P_0)$ of automorphisms fixing the origin has order $4$, so by (\Goto{result-IV-4.7}{4.7}) we must have $j=1728$}

\ChapterEntry{p-346}{IV, Hasse invariant. Semilinearity of Frobenius}{332}


\Printed{This map is not linear, but it is \emph{$p$-linear}, namely $F^*(\lambda a)=\lambda^p F^*(a)$ for all $\lambda\in k$, $a\in H^1(X,\mathcal O_X)$.}
\Revised{This map is $p$-semilinear, namely $F^*(\lambda a)=\lambda^p F^*(a)$ for all $\lambda\in k$, $a\in H^1(X,\mathcal O_X)$. If it is nonzero and Frobenius on $k$ is nontrivial, it is not $k$-linear.}

\ChapterEntry{p-347}{IV, Corollary 4.22. Sign of the Legendre coefficient}{333}


\Printed{Summing up, the coefficient of $(xyz)^{p-1}$ in $f^{p-1}$ is
\[
(-1)^k\binom{p-1}{k}\sum\binom{k}{i}^{\!2}\lambda^i.
\]
Since the outer factor is $\equiv1\pmod p$, we get $h_p(\lambda)$ as defined above.}
\Revised{Each product of those two coefficients contributes a sign $(-1)^k$, which cancels the sign $(-1)^k$ from choosing $k$ copies of $-x(x-z)(x-\lambda z)$. Thus the coefficient of $(xyz)^{p-1}$ in $f^{p-1}$ is
\[
\binom{p-1}{k}\sum\binom{k}{i}^{\!2}\lambda^i.
\]
In characteristic $p$ this is $(-1)^k h_p(\lambda)$, so the vanishing criterion is as stated.}

\Printed{Since the correspondence $\lambda\mapsto j$ is $6$ to $1$ with two exceptions, we can have at most $k/6+2$ values of $j$, hence at most $\lfloor p/12\rfloor+2$.}
\Revised{For $p>3$, the correspondence $\lambda\mapsto j$ is $6$ to $1$ away from two exceptions, so there are at most $k/6+2$ values of $j$, hence at most $\lfloor p/12\rfloor+2$. If $p=3$, then $h_3(\lambda)=1+\lambda$. Its only root is $-1$, which gives $j=0=1728$ and satisfies the same bound.}

\ChapterEntry{p-348}{IV, Corollary 4.23 and following note. Characteristic three}{334}


\Printed{$j=0$ occurs $\Leftrightarrow p\equiv2\pmod3$ (Ex.~\Goto{ex-IV-4.14}{4.14}}
\Revised{for $p>3$, $j=0$ occurs $\Leftrightarrow p\equiv2\pmod3$ (Ex.~\Goto{ex-IV-4.14}{4.14}}

\Printed{There are also tables of these $j$ for small values of $p$---see Deuring}
\Revised{For $p=3$ the two exceptional values coincide and $j=0=1728$ is the unique supersingular value. There are also tables of these $j$ for small values of $p$---see Deuring}

\ChapterEntry{p-348}{IV, Example 4.23.2. Choice of the cube root of unity}{334}


\Printed{$\omega^3=1$. So $j=0$.}
\Revised{$\omega$ a primitive cube root of unity. So $j=0$.}

\ChapterEntry{p-350}{IV, Exercise 4.4. The invariant in positive characteristic}{336}


\Printed{Show that the $j$-invariant is a rational function of the $a_i$, with coefficients in $\mathbb Q$.}
\Revised{Show that $j=c_4^3/\Delta$, with $c_4,\Delta$ the integer-coefficient Weierstrass invariants and $\Delta\ne0$.}

\ChapterEntry{p-351}{IV, Exercise 4.5. Characteristic for the three degree-two cases}{337}


\Printed{Let $X,P_0$ be an elliptic curve having an endomorphism $f\colon X\to X$}
\Revised{Assume $\operatorname{char}k=0$. Let $X,P_0$ be an elliptic curve having an endomorphism $f\colon X\to X$}

\ChapterEntry{p-351}{IV, Exercise 4.6. Weights of inflection and hyperosculation}{337}


\Printed{(a) Let $X$ be a curve of genus $g$ embedded birationally in $\mathbf P^2$}
\Revised{Assume $\operatorname{char}k=0$ throughout. (a) Let $X$ be a curve of genus $g$ embedded birationally in $\mathbf P^2$}

\Printed{$X$ has $6(g-1)+3d$ inflection points. A node does not count as an inflection point. Assume $\operatorname{char}k=0$.}
\Revised{the inflection divisor on the normalization has degree $6(g-1)+3d$, including branches at nodes with positive inflection weight.}

\Printed{induction on $n$, to show that $X$ has $n(n+1)(g-1)+(n+1)d$ hyperosculation points}
\Revised{induction on $n$, to show that $X$ has $n(n+1)(g-1)+(n+1)d$ hyperosculation points counted with ramification weights $\sum_{i=0}^n(a_i-i)$ for vanishing sequence $(a_i)$}

\Printed{points of order $d$ in its group law.}
\Revised{points of order dividing $d$ in its group law.}

\ChapterEntry{p-351}{IV, Exercise 4.7(d). Degree zero in the additivity argument}{337--338}

\PassagePage{337}


\Printed{$\mathcal L\in\operatorname{Pic}X'$}
\Revised{$\mathcal L\in\operatorname{Pic}^{0}X'$}

\PassagePage{338}


\Printed{on $X'$, $p_2^*\mathcal L\in\operatorname{Pic}_\sigma^0$ to prove the result.]}
\Revised{in $\operatorname{Pic}^0X'$, $p_2^*\mathcal L\in\operatorname{Pic}_\sigma^0$ to prove the result.]}

\ChapterEntry{p-352}{IV, Exercise 4.7(c), (e), (f). Degree of finite maps}{337--338}

\PassagePage{337}


\Printed{Assume $f(P_0)=P'_0$, and let $n=\deg f$.}
\Revised{Assume $f(P_0)=P'_0$ and $f$ is nonconstant, and let $n=\deg f$.}

\PassagePage{338}


\Printed{$n\in\mathbb Z$, $\widehat n_X=n_X$. Conclude that $\deg n_X=n^2$}
\Revised{$n\in\mathbb Z$, $\widehat n_X=n_X$. Conclude that $\deg n_X=n^2$ for $n\ne0$}

\Printed{$f$ that $\deg\widehat f=\deg f$.}
\Revised{nonconstant $f$ that $\deg\widehat f=\deg f$.}

\ChapterEntry{p-352}{IV, Exercise 4.8. Product notation for the fundamental group}{338}

\Printed{\[
\begin{aligned}
\pi_1(X)&=\prod_{\ell\text{ prime}}\mathbb Z_\ell\times\mathbb Z_\ell,\\
\pi_1(X)&=\prod_{\ell\ne p}\mathbb Z_\ell\times\mathbb Z_\ell,\\
\pi_1(X)&=\mathbb Z_p\times\prod_{\ell\ne p}\mathbb Z_\ell\times\mathbb Z_\ell.
\end{aligned}
\]}
\Revised{\[
\begin{aligned}
\pi_1(X)&\cong\prod_{\ell\text{ prime}}(\mathbb Z_\ell\times\mathbb Z_\ell),\\
\pi_1(X)&\cong\prod_{\ell\ne p}(\mathbb Z_\ell\times\mathbb Z_\ell),\\
\pi_1(X)&\cong\mathbb Z_p\times\prod_{\ell\ne p}(\mathbb Z_\ell\times\mathbb Z_\ell).
\end{aligned}
\]
In the products, $\ell$ runs over the indicated primes.}

\ChapterEntry{p-352}{IV, Exercise 4.9(b). Kernel of an inseparable isogeny}{338}


\Printed{$X'$ is uniquely determined by $X$ and $\ker f$.}
\Revised{Translate the target to make $f$ pointed. Then $X'$ is determined by $X$ and the scheme-theoretic kernel of $f$. Factor out the purely inseparable part by a Frobenius twist and count the finite subgroups of geometric points for the separable part.}

\ChapterEntry{p-353}{IV, Exercise 4.15. The subgroup killed by $p$}{339}


\Printed{points of order $p$ on $X$ is isomorphic to $\mathbb Z/p$;}
\Revised{points of order dividing $p$ on $X$ is isomorphic to $\mathbb Z/p$;}

\ChapterEntry{p-353}{IV, Exercise 4.16. Rational origin and the degree inequality}{339}


\Printed{Assume also that $X$ has a rational point over $\mathbb F_q$.}
\Revised{Choose an $\mathbb F_q$-rational point as the origin of $X$.}

\Printed{Use the fact that $\deg(m+nF')>0$}
\Revised{Set the degree of the zero endomorphism to $0$. Use the fact that $\deg(m+nF')\ge0$}

\ChapterEntry{p-354}{IV, Exercise 4.19. The plane model and the multiplication map}{340}


\Printed{an equation with integer coefficients}
\Revised{a primitive integral homogeneous cubic}

\Printed{consisting of all primes $p\ne2$ such that the fibre $X_{(p)}$ of $X$ over $p$ is nonsingular. For any $n$}
\Revised{consisting of the generic point and the primes $p\ne2$ at which the plane cubic fibre $X_{(p)}$ is smooth. The origin extends to a section over $T$. For any $n\ge1$}

\Printed{$p\in T$, the natural map $X(\mathbb Q)\to X_{(p)}(\mathbb F_p)$ induced on}
\Revised{closed point $p\in T$, the natural map $X(\mathbb Q)\to X_{(p)}(\mathbb F_p)$ induced on}

\ChapterEntry{p-354}{IV, Exercise 4.20(a)--(b). The Frobenius twist and the ground field}{340}


\Printed{ring of endomorphisms.}
\Revised{ring of endomorphisms. For parts (a) and (b), take $k$ algebraically closed.}

\Printed{If Hasse $(X)=0$ show in any case $j\in\mathbb F_{p^2}$}
\Revised{If Hasse $(X)=0$ show in any case $j\in\mathbb F_{p^2}$, which can be checked after extending $k$ to an algebraic closure}

\ChapterEntry{p-354}{IV, Exercise 4.22. Trivializing the family after fixing the cross-ratio}{340}


\Printed{, so the family is trivial.}
\Revised{. Use three branch sections to trivialize the quotient projective-line bundle and the constant cross-ratio to fix the fourth. Then show that the remaining quadratic twist is trivial: the square root of the branch line bundle is unique, and every unit of the base is a square.}

\ChapterEntry{p-357}{IV, Proposition 5.3. Scheme fibers in the canonical morphism}{343}


\Printed{Since the $g_2^1$ has infinitely many divisors in it, we see that $f$ cannot be birational.}
\Revised{The degree-two map $f_0$ of this pencil is separable. In characteristic two a purely inseparable quadratic extension of $k(t)$ is $k(t^{1/2})$, which would make $X$ rational. Thus a general fibre has two distinct points, and $f$ identifies them, so $f$ cannot be birational.}

\Printed{Since $f$ already collapses the pairs of the $g_2^1$ we chose above, it must be equal to the composition of the map $f_0\colon X\to\mathbf P^1$ determined by our $g_2^1$ with the $(g-1)$-uple embedding of $\mathbf P^1$.}
\Revised{Since $f$ identifies the two points of a general $f_0$-fibre, its pullback of rational functions lies in the fixed field of the involution of $f_0$. Hence $f$ factors through $f_0$. Both morphisms have degree two, so the induced map from $\mathbf P^1$ to $X'$ is an isomorphism, followed by the $(g-1)$-uple embedding in $\mathbf P^{g-1}$.}

\Printed{set of $g-1$ points of $X'$ is a hyperplane section, so we can identify the canonical}
\Revised{effective divisor of degree $g-1$, with multiplicities, of $X'$ is a hyperplane section, so we can identify the canonical}

\ChapterEntry{p-357}{IV, Theorem 5.4. Linear equivalence in Clifford's equality case}{343--345, 351}

\PassagePage{343}


\Printed{either $D=0$ or $D=K$ or $X$ is hyperelliptic and $D$ is a multiple of the unique $g_2^1$}
\Revised{either $D=0$ or $D\sim K$ or $X$ is hyperelliptic and $D$ is a multiple of the unique $g_2^1$}

\PassagePage{344}


\Printed{then $K-D$ is also effective, so we can apply the lemma, and obtain}
\Revised{then $K-D$ has an effective representative, so we can apply the lemma and obtain}

\Printed{case $D=0$ or $D=K$}
\Revised{case $D=0$ or $D\sim K$}

\Printed{$D\ne0,K$, and that $\dim|D|=\frac12\deg D$.}
\Revised{$D\ne0$ and $D\not\sim K$, and that $\dim|D|=\frac12\deg D$.}

\PassagePage{345}


\Printed{$D\ne0,K$, and $\dim|D|=\frac12\deg D$.}
\Revised{$D\ne0$ and $D\not\sim K$, and $\dim|D|=\frac12\deg D$.}

\Needspace{3\baselineskip}
\PassagePage{351}


\Printed{is absurd) or $D=K$, in which case $X$ is the canonical curve of genus $4$, or $X$ is}
\Revised{is absurd) or $D\sim K$, in which case $X$ is the canonical curve of genus $4$, or $X$ is}

\ChapterEntry{p-359}{IV, Theorem 5.4. Maps in the intersection-divisor sequence}{345}


\Printed{the first map is addition, the second subtraction.}
\Revised{the first map is the diagonal $s\mapsto(s,s)$, and the second is subtraction $(u,v)\mapsto u-v$.}

\ChapterEntry{p-360}{IV, Example 5.5.3. Trisecants and divisors with repeated points}{346}


\Printed{we show that $X$ has a $g_3^1$ if and only if it has a trisecant in this embedding. Let $P,Q,R\in X$.}
\Revised{we show that $X$ has a $g_3^1$ if and only if it has a trisecant in this embedding, meaning a line containing a length-three divisor of $X$. Let $P,Q,R\in X$.}

\Printed{$P,Q,R$. Hence $\dim|P+Q+R|=1$ if and only if $P,Q,R$ are contained in a $2$-dimensional family of hyperplanes, which is equivalent to saying that $P,Q$ and $R$ are collinear.}
\Revised{the divisor $P+Q+R$ as a subscheme, including its multiplicities. Hence $\dim|P+Q+R|=1$ if and only if this subscheme is contained in a $2$-dimensional family of hyperplanes, or equivalently, in a line.}

\Printed{$L$ would meet each of the quadric hypersurfaces in three points, so it would have to be contained in these hypersurfaces,}
\Revised{each quadratic equation on $L$ would vanish on a subscheme of length three, so $L$ would be contained in all three hypersurfaces,}

\Printed{So let $Q,R,S$ lie on a trisecant of $X'$. Then their inverse images on $X$, together with $P$, form four points which lie in a plane of $\mathbf P^4$. Then the same argument as above shows these points give a $g_4^1$.}
\Revised{Take a length-three divisor on $X'$ contained in a line. In the linear system $|K-P|$, its inverse image lies in all hyperplanes through that line. Its inverse image plus $P$ is a degree-four divisor in a plane of $\mathbf P^4$, with multiplicities included. The same dimension argument gives a $g_4^1$.}

\ChapterEntry{p-361}{IV, Example 5.5.5. Branch-point description in characteristic two}{347}


\Printed{using the method of (Ex.~\Goto{ex-IV-2.2}{2.2}), one can show that}
\Revised{when $\operatorname{char}k\ne2$, using the method of (Ex.~\Goto{ex-IV-2.2}{2.2}), one can show that}

\Printed{of dimension $3$.}
\Revised{of dimension $3$. In characteristic two the branch description requires wild ramification data.}

\ChapterEntry{p-362}{IV, Exercise 5.5(b)--(c). The singular point of the plane quintic}{348}


\Printed{one node. These form an irreducible family of dimension $11$.}
\Revised{one double point of $\delta$-invariant one, either a node or a cusp. These form an irreducible family of dimension $11$.}

\Printed{conics through the node cut out the canonical system (not counting the fixed points at the}
\Revised{conics through the double point cut out the canonical system (not counting the fixed points at the}

\Printed{node). Mapping $\mathbf P^2\to\mathbf P^4$ by this linear system of conics,}
\Revised{double point). Mapping $\mathbf P^2\to\mathbf P^4$ by this linear system of conics,}

\ChapterEntry{p-362}{IV, Exercise 5.7(a). Existence of the canonical embedding}{348}


\Printed{curve of genus $3$ is induced by an automorphism of $\mathbf P^2$ via}
\Revised{nonhyperelliptic curve of genus $3$ is induced by an automorphism of $\mathbf P^2$ via}

\ChapterEntry{p-362}{IV, Exercise 5.7(b)--(c). Characteristic restrictions for automorphisms}{348--349}

\PassagePage{348}


\Printed{Assume $\operatorname{char}k\ne3$. If $X$ is the curve given by}
\Revised{Assume $\operatorname{char}k\ne3$ and $\operatorname{char}k\ne7$. If $X$ is the curve given by}

\PassagePage{349}


\Printed{. See Burnside \Goto{bib-474-8-item-1}{[1, \S232]} or Klein}
\Revised{in characteristic zero. See Burnside \Goto{bib-474-8-item-1}{[1, \S232]} or Klein}

\Printed{For example, if $n=2$, then for suitable choice of coordinates, $T$ can be written as}
\Revised{For example, if $n=2$ and $\operatorname{char}k\ne2$, then for suitable choice of coordinates, $T$ can be written as}

\ChapterEntry{p-363}{IV, Exercise 5.7(c). Order of an involution}{349}


\Printed{an automorphism of degree $2$ form a family of dimension $12$ inside the $14$-dimensional family}
\Revised{an automorphism of order $2$ form a family of dimension $12$ inside the $14$-dimensional family}

\ChapterEntry{p-364}{IV, Proposition 6.1. Dimension increase under linear equivalence}{350}


\Printed{these two is $1$, so we see that the set of $D\in X^d$ which are linearly equivalent}
\Revised{these two is at most $1$, so we see that the set of $D\in X^d$ which are linearly equivalent}

\ChapterEntry{p-369}{IV, Exercise 6.8. The zero divisor on a rational curve}{355}


\Printed{A curve $X$ of genus $g$ has a nonspecial divisor $D$ of degree $d$ such that $|D|$ has no base points if and only if $d\ge g+1$.}
\Revised{A curve $X$ of genus $g\ge1$ has a nonspecial divisor $D$ of degree $d$ such that $|D|$ has no base points if and only if $d\ge g+1$. For $g=0$, the condition is $d\ge0$.}

\ChapterEntry{p-369}{IV, Exercise 6.9. Smoothness along the prescribed curve}{355}


\Printed{Apply Bertini's theorem to the projective embedding of $\widetilde{\mathbf P}$ corresponding to $\mathcal I_Y\otimes\pi^*\mathcal O_{\mathbf P^3}(m)$.}
\Revised{Apply Bertini's theorem to the projective embedding of $\widetilde{\mathbf P}$ corresponding to $\mathcal I_Y\otimes\pi^*\mathcal O_{\mathbf P^3}(m)$. For smoothness along $X$, also choose the defining section so that its image in $(\mathcal I_X/\mathcal I_X^2)(m)$ is nowhere zero. The latter is a globally generated rank-two bundle for large $m$, so apply (II, Ex.~\Goto{ex-II-8.2}{8.2}) and lift the chosen section using $H^1(\mathcal I_X^2(m))=0$.}

